%=====================================================================
\section{Limitations and next steps}
\label{sec:limits}

\subsection{Limitations}

\begin{enumerate}
  \item \textbf{No optimality guarantee.} The policy is a heuristic. Its plans are judged only by comparison with the MILP; on instances the MILP cannot solve to optimality, the true gap of both is unknown.
  \item \textbf{Hard constraints can still defeat it.} Masks guarantee that every \emph{placed} EV is feasible, but a sequence of locally reasonable choices can leave a later EV without options. Such plans are reported as invalid (58 of 560 for RL1, 51 of them proven infeasible).
  \item \textbf{Training distribution.} The policy is only as good as the generator's resemblance to the test files (Assumption~\ref{as:dist}). It was trained on the templates of the 100-EV files; its quality on the 200-EV files (more stations and vans) is not measured.
  \item \textbf{Teacher bias.} The warm-up imitates the greedy rule; REINFORCE then improves on it, but some of the rule's habits may survive. The validation yardstick is also the rule, not the MILP.
  \item \textbf{Hyperparameters not tuned.} Network size follows Kool et al.; learning rates, batch sizes, penalty, and the active-search settings ($\lambda$, learning rate, samples, time cap) were set by hand. The EAS paper tunes its settings by Bayesian optimisation; this project did not.
  \item \textbf{One training run.} RL1's numbers come from one training seed. Differences between RL1 and RL2 smaller than the run-to-run variation of training cannot be attributed to the changes.
  \item \textbf{Small evidence for active search.} The active-search gain was measured on three instances on a CPU with 30 iterations. The server run on all 560 instances is the real test.
  \item \textbf{Static problem only.} Like formulation 3, the policy assumes all requests are known at time 0. It does not yet handle requests that arrive during the period.
  \item \textbf{Compute.} Training needs a GPU for a reasonable time (39 minutes on an RTX 4090 for RL1; hours on a laptop CPU). Solving is cheap on any machine.
\end{enumerate}

\subsection{Next steps}

\begin{enumerate}
  \item \textbf{Run RL2 on the server} and its ablation on RL1's model (\code{-{}-run-id-prefix rl1model}) to separate the effect of training from that of active search.
  \item \textbf{Several training seeds} for RL1 and RL2, to put confidence intervals on the gaps.
  \item \textbf{Tune active search} ($\lambda$, learning rate, samples per iteration, time cap) on separate validation instances, and try the other two EAS variants (an added residual layer, EAS-Lay; a lookup table, EAS-Tab).
  \item \textbf{RL plans as MILP start solutions.} RL1 finds valid plans for 23 instances in which the MILP finds none in 300 s; handing them to Gurobi as a start (as the greedy plan already is) could help the exact solver, without restricting its model.
  \item \textbf{A learned improvement stage} (neural local search) applied to RL plans, which are feasible starting points that greedy cannot provide for the hard R1/RC1 instances.
  \item \textbf{Training on the 200-EV templates} (\code{-{}-file-sizes 100 200}) and testing on the Gehring--Homberger files.
  \item \textbf{An online version.} If EVs are processed in request order and the decoder only chooses the resource, the same network becomes a dispatcher for requests that arrive during the period (the dynamic extension of the project brief).
\end{enumerate}
